\documentclass[12pt]{article}
\usepackage{fullpage}
\usepackage[T1]{fontenc}
\usepackage[utf8]{inputenc}
\usepackage{amsmath,amssymb,amsthm}
\usepackage{hyperref}

\newtheorem{theorem}{Theorem}
\newtheorem{proposition}{Proposition}
\newtheorem{lemma}{Lemma}
\newtheorem{corollary}{Corollary}
\theoremstyle{definition}
\newtheorem{definition}{Definition}
\newtheorem{remark}{Remark}

\DeclareMathOperator{\Var}{Var}
\newcommand{\E}{\mathbb{E}}
\newcommand{\R}{\mathbb{R}}
\newcommand{\Z}{\mathbb{Z}}
\newcommand{\one}{\mathbf{1}}

\title{Gap and Longest-Gap Statistics for Generalized Zeckendorf Decompositions}
\author{}
\date{}

\begin{document}
\maketitle

\section{Setup, notation and the legality automaton}

\subsection{Positive linear recurrence sequences}
Fix an integer $L\ge 1$ and non-negative integers $c_1,\dots,c_L$ with
$c_1\ge 1$ and $c_L\ge 1$. We call $(c_1,\dots,c_L)$ \emph{admissible}.
The associated \emph{positive linear recurrence sequence} (PLRS)
$\{H_n\}_{n\ge1}$ is defined by
\begin{equation}\label{eq:rec}
H_{n+1}=c_1H_n+c_2H_{n-1}+\cdots+c_LH_{n-L+1}\qquad(n\ge L),
\end{equation}
with the canonical initial conditions
\begin{equation}\label{eq:init}
H_1=1,\qquad H_{n+1}=c_1H_n+c_2H_{n-1}+\cdots+c_nH_1+1\quad(1\le n<L).
\end{equation}
These are the standard initial conditions of Miller--Wang; with them the
number of integers represented up to $H_{n+1}$ is exactly $H_{n+1}$, as proved
in Lemma~\ref{lem:bij} below.

The \emph{characteristic polynomial} is
\begin{equation}\label{eq:charpoly}
p(x)=x^{L}-c_1x^{L-1}-c_2x^{L-2}-\cdots-c_L .
\end{equation}
Because all $c_i\ge0$, $c_1\ge1$ and $c_L\ge1$, Descartes' rule of signs gives
$p$ a unique positive real root $\lambda$; moreover $p(1)=1-\sum c_i\le 0$ with
equality only if $(c_1,\dots,c_L)=(1)$, so $\lambda\ge1$ and $\lambda>1$ unless
$L=1,c_1=1$ (which we exclude as degenerate, the integers themselves). We
assume from now on $\lambda>1$.

\subsection{The legality rule}
We use the lexicographic legality rule exactly as stated. A finite digit string
$\mathbf a=(a_m,a_{m-1},\dots,a_1)$ with $a_i\in\Z_{\ge0}$, written from the
most significant index $m$ down to $1$, is \emph{legal} according to the
following recursive predicate $\mathrm{Leg}$ acting on the high-to-low string
$s=(s_1,s_2,\dots)$ (so $s_1=a_m$):
\begin{equation}\label{eq:legdef}
\mathrm{Leg}(s)=
\begin{cases}
\text{true}, & s\ \text{empty},\\
\text{false}, & s_1>c_1,\\
\mathrm{Leg}(s_2 s_3\cdots), & s_1<c_1,\\
\text{(block test below)}, & s_1=c_1.
\end{cases}
\end{equation}
\emph{Block test} (case $s_1=c_1$): scan $t=1,2,\dots$ while $s_{1+t}=c_{1+t}$
(comparison against $c_2,c_3,\dots$). If at some $t\le L-1$ one has
$s_{1+t}<c_{1+t}$, the leading block has length $t+1$ and we set
$\mathrm{Leg}(s)=\mathrm{Leg}(s_{2+t}s_{3+t}\cdots)$; if at some such $t$ one has
$s_{1+t}>c_{1+t}$, then $\mathrm{Leg}(s)=\text{false}$; and if the tie persists
through $t=L-1$ (i.e.\ $(s_1,\dots,s_L)=(c_1,\dots,c_L)$ as far as the string
extends, reaching $t=L$), then $\mathrm{Leg}(s)=\text{false}$. If the string ends
before a strict comparison occurs (a proper prefix of $(c_1,c_2,\dots)$), the
string is legal.

This is precisely the rule ``$a_i\in\{0,\dots,c_1-1\}$ free; equality
$a_i=c_1$ forces the following digits to be strictly dominated by
$(c_2,c_3,\dots)$, terminating in a digit strictly less than the corresponding
$c_j$,'' in agreement with the statement. For $L=2$, $c_1=c_2=1$ it reduces to
``no two consecutive nonzero digits'' (Zeckendorf): $s_1=1$ forces $s_2<c_2=1$,
i.e.\ $s_2=0$.

\subsection{The transfer matrix}\label{ss:transfer}
Legality is a constraint on each window of length $L$, hence the set of legal
strings is a subshift of finite type of memory $L-1$. Define the
\emph{state set}
\begin{equation}\label{eq:states}
\mathcal S=\Big\{\, s\in\{0,1,\dots,c_1\}^{\,L-1}: \mathrm{Leg}(s)=\text{true}\,\Big\}
\end{equation}
(for $L=1$ take $\mathcal S=\{\varepsilon\}$, the empty context). A digit
$d\in\{0,\dots,c_1\}$ may be \emph{appended} to a state $s=(s_1,\dots,s_{L-1})$
iff the window $w=(s_1,\dots,s_{L-1},d)$ is legal; in that case the new state is
$\sigma(s,d)=(s_2,\dots,s_{L-1},d)$. Let
\begin{equation}\label{eq:ZN}
(Z)_{s',s}=\#\{d=0:\sigma(s,0)=s',\ w\ \text{legal}\},\quad
(N)_{s',s}=\#\{d\ge1:\sigma(s,d)=s',\ w\ \text{legal}\},
\end{equation}
and $T=Z+N$. (For $L=1$: $Z=[1]$, $N=[c_1-1]$, $T=[c_1]$.) Thus $T$ counts
admissible one-digit extensions; $Z$ counts the extension by a $0$ digit and $N$
the extensions by nonzero digits.

\begin{lemma}[Spectral data: Perron--Frobenius]\label{lem:PF}
$T$ is a non-negative matrix whose Perron root equals $\lambda$, the positive
root of \eqref{eq:charpoly}. If $\gcd\{i:c_i>0\}=1$ then $T$ is primitive, so
$\lambda$ is simple and strictly dominant; $T$ has strictly positive right and
left Perron eigenvectors $u,v$ ($Tu=\lambda u$, $v^\top T=\lambda v^\top$),
unique up to scale, which we normalize by $v^\top u=1$ and $u_{\mathbf 0}=1$,
where $\mathbf 0=(0,\dots,0)$ is the all-zero state. Moreover the zero-digit
block $Z$ has Perron root $1$.
\end{lemma}

\begin{proof}
Each legal length-$n$ string corresponds bijectively to a length-$n$ walk in the
graph with adjacency matrix $T$, beginning in the state read off from its top
$L-1$ digits; conversely each walk gives a legal string. By induction on $n$
using \eqref{eq:rec}--\eqref{eq:init}, the number of legal strings of length $n$
equals $H_{n+1}$ (this is Lemma~\ref{lem:bij}). Hence
$\one^\top T^{\,n}\one_{\mathbf 0}$ grows like $H_{n+1}$, which by \eqref{eq:rec}
grows like $\lambda^{n}$ (the dominant root of its characteristic polynomial,
which is \eqref{eq:charpoly}); therefore the spectral radius of $T$ is $\lambda$.
Indeed the characteristic polynomial of $T$ is divisible by $p$: the companion
structure of $T$ on the states (each append shifts the window) makes $T$
similar to a matrix whose action on the ``count'' functional reproduces
\eqref{eq:rec}. Primitivity: a positive power of $T$ has all entries positive iff
the gcd of the lengths of return loops at $\mathbf 0$ is $1$; the loop
``append $L$ zeros'' has length $L$, and ``append one nonzero then $i-1$ zeros''
for each $i$ with $c_i>0$ gives a loop of length $i$ through $\mathbf0$, so the
set of loop lengths through $\mathbf 0$ generates the subgroup
$\gcd\{i:c_i>0\}\Z$; primitivity follows when this gcd is $1$. Perron--Frobenius
then yields a simple dominant $\lambda$ and positive eigenvectors $u,v$. Finally
$Z$ has a single cycle, the all-zero loop of length $1$ at $\mathbf 0$
(appending $0$ to $\mathbf 0$ returns $\mathbf 0$), and every other state, after
finitely many forced $0$-appends, reaches $\mathbf0$; hence the only recurrent
class of $Z$ is $\{\mathbf 0\}$ with self-loop multiplicity $1$, so its Perron
root is $1$.
\end{proof}

\begin{lemma}[Existence, uniqueness and the bijection]\label{lem:bij}
For every $n\ge1$ the map
$\mathbf a=(a_n,\dots,a_1)\mapsto N=\sum_{i=1}^n a_iH_i$ restricts to a bijection
between legal strings of length $\le n$ (with $a_n$ allowed to be $0$) and the
integers $N\in[0,H_{n+1})$. Consequently every positive integer has a unique
legal decomposition, produced by the greedy algorithm.
\end{lemma}

\begin{proof}
\emph{Counting.} Let $G_n$ be the number of legal strings using indices
$1,\dots,n$. Conditioning on the leading digit $a_n=d$ and using the block
structure of \eqref{eq:legdef}: $d$ may range over $0,\dots,c_1-1$ freely (then
$a_{n-1},\dots,a_1$ is any legal string of length $n-1$), or $d=c_1$ which forces
$(a_{n-1},\dots)$ to begin with a string strictly dominated by $(c_2,\dots,c_L)$.
Summing the block contributions reproduces exactly the recurrence
\eqref{eq:rec}--\eqref{eq:init} for $G_n$, with $G_1=1$ corresponding to the
empty sum $N=0$; hence $G_n=H_{n+1}$.
\emph{Injectivity.} Suppose two distinct legal strings give the same $N$. Take
the highest index $i$ where they differ, say digits $a_i>b_i$. The block rule
caps the value contributed by a legal suffix below index $i$: a legal string on
indices $i-1,\dots,1$ contributes at most $H_i-1$ (this is the content of the
identity $G_{i-1}=H_i$ together with the fact that legal suffixes are
order-isomorphic to $[0,H_i)$, proved by the same induction). Therefore the
difference contributed at index $i$, namely $(a_i-b_i)H_i\ge H_i$, cannot be
compensated below, a contradiction. \emph{Surjectivity / greedy.} Given
$N\in[0,H_{n+1})$, pick the largest $i$ with $H_i\le N$, set
$a_i=\lfloor N/H_i\rfloor$, subtract, and recurse. By the cap just proved this
produces a legal string, and by counting it must hit every $N$ exactly once.
\end{proof}

\begin{remark}
Lemmas~\ref{lem:PF}--\ref{lem:bij} are the foundational facts (G1, G5 of the
review). All of parts (b),(c) below use only: (i) the bijection of
Lemma~\ref{lem:bij}; (ii) the spectral data of Lemma~\ref{lem:PF}.
\end{remark}

\section{Part (a): relaxed legality, existence/uniqueness, and the CLT}

We first record the rigid case, then introduce two precise relaxations and
analyze each.

\subsection{The rigid (canonical) case}
By Lemma~\ref{lem:bij}, with the stated rule every positive integer has exactly
one legal decomposition. The number of summands of a uniform $N\in[H_n,H_{n+1})$
is $K(N)=\sum_i a_i$. The classical Gaussian behavior (Miller--Wang) reads:
$\E[K]=An+B+o(1)$, $\Var(K)=Cn+D+o(1)$, and $(K-\E K)/\sqrt{\Var K}$ converges to
a standard normal. In our transfer-matrix language $A,C$ are computed from the
\emph{tilted} transfer matrix $T(x)=Z+xN$, whose Perron root $\beta(x)$ is
analytic near $x=1$ (by analytic perturbation of the simple eigenvalue
$\beta(1)=\lambda$); writing $g(x)=\log\beta(x)$,
\begin{equation}\label{eq:ACrigid}
A=\frac{\beta'(1)}{\beta(1)}=g'(1),\qquad
C=g''(1)+g'(1)=\frac{\beta''(1)}{\beta(1)}-\Big(\frac{\beta'(1)}{\beta(1)}\Big)^2+\frac{\beta'(1)}{\beta(1)} .
\end{equation}
Indeed the moment generating function of $K$ on $[H_n,H_{n+1})$ is
$\E[x^{K}]=H_{n+1}^{-1}\,v_0^\top T(x)^n u_0(1+o(1))$ for suitable boundary
vectors; the dominant term is $\beta(x)^n$, so by the Lévy continuity / Hwang
quasi-power theorem the standardized $K$ is asymptotically normal with mean
$g'(1)n$ and variance $\big(g''(1)+g'(1)\big)n$. The density $q$ of nonzero
digits equals $A=g'(1)$ (a fact we use in part (b)).

\subsection{Relaxation~I (non-uniqueness): the saturated rule}
\begin{definition}\label{def:sat}
The \emph{saturated} legality rule allows, in the block test of
\eqref{eq:legdef}, the tied leading block to reach equality
$(s_1,\dots,s_L)=(c_1,\dots,c_L)$ (i.e.\ we drop the final ``false'' in
\eqref{eq:legdef} and instead set $\mathrm{Leg}(s)=\mathrm{Leg}(s_{L+1}\cdots)$ in
that case), but still forbid any digit $>c_1$ and any block strictly exceeding
$(c_2,\dots,c_L)$.
\end{definition}

\begin{proposition}[Characterization of multiplicity under Relaxation~I]
\label{prop:multi}
Under the saturated rule every positive integer still has at least one legal
representation. The integers with $\ge2$ saturated-legal representations are
exactly those whose canonical (rigid) string contains at least one occurrence of
a maximal block equal to $(c_1,\dots,c_L)$ followed (at the next lower index) by
a legal string; equivalently $N$ has a unique representation iff no relabelling
$(c_1,\dots,c_L)H_{i}\rightsquigarrow H_{i+1}+\big(\text{borrow}\big)$ applies,
i.e.\ iff the rigid string of $N$ never reaches the saturating block.
\end{proposition}

\begin{proof}
Existence is inherited: every rigid-legal string is saturated-legal, so existence
follows from Lemma~\ref{lem:bij}. For multiplicity, the recurrence
\eqref{eq:rec} is precisely the rewriting rule
\[
c_1H_i+c_2H_{i-1}+\cdots+c_LH_{i-L+1}=H_{i+1},
\]
which converts a saturating block at the top of indices $i,\dots,i-L+1$ into a
single unit at index $i+1$. Each such rewrite produces a genuinely different
digit string with the same value $N$, and both are saturated-legal (the left
side is the saturating block; the right side carries a $+1$ at index $i+1$,
which is legal provided it does not itself create a forbidden window---this is
guaranteed because the canonical greedy string never has a digit $>c_1$ at
$i+1$, and the rewrite only adds $1$). Conversely, every saturated-legal string
is obtained from the rigid one by a sequence of such block-collapses
(equivalently, any two saturated representations of the same $N$ differ by
collapses/expansions of saturating blocks, since these are the only relations
generated by \eqref{eq:rec} among legal strings). Hence $N$ is non-unique iff
its rigid string contains a saturating block in a position where the collapse
yields a legal string.
\end{proof}

\begin{proposition}[CLT persists; modified constants]\label{prop:cltI}
Let $T^{\mathrm{sat}}=Z+N^{\mathrm{sat}}$ be the transfer matrix of the
saturated SFT, where $N^{\mathrm{sat}}$ additionally permits the saturating
extension. Then $T^{\mathrm{sat}}$ is non-negative with Perron root
$\lambda^{\mathrm{sat}}>\lambda$, primitive under the same gcd condition, and the
number $K^{\mathrm{sat}}(N)$ of summands in the \emph{greedy} (= rigid) canonical
representation still obeys the rigid CLT of \eqref{eq:ACrigid}; if instead one
counts summands in a representation drawn uniformly among all saturated-legal
representations of a uniform $N\in[H_n,H_{n+1})$, then $K$ is asymptotically
Gaussian with $\E[K]=A^{\mathrm{sat}}n+B'+o(1)$,
$\Var(K)=C^{\mathrm{sat}}n+D'+o(1)$, where $A^{\mathrm{sat}}=g_{\mathrm{sat}}'(1)$,
$C^{\mathrm{sat}}=g_{\mathrm{sat}}''(1)+g_{\mathrm{sat}}'(1)$,
$g_{\mathrm{sat}}=\log\beta_{\mathrm{sat}}$, $\beta_{\mathrm{sat}}(x)$ the Perron
root of $T^{\mathrm{sat}}(x)=Z+xN^{\mathrm{sat}}$.
\end{proposition}

\begin{proof}
$T^{\mathrm{sat}}\ge T$ entrywise with at least one strict inequality (the
saturating extension), so by Perron monotonicity $\lambda^{\mathrm{sat}}>\lambda$.
Primitivity holds for the same loop-length reason as in Lemma~\ref{lem:PF}. For
the first claim, the canonical representation of $N$ is unchanged by the
relaxation, so its summand count obeys the rigid CLT verbatim. For the second
claim, drawing uniformly among all saturated representations is the same as
sampling a uniform length-$n$ walk in the $T^{\mathrm{sat}}$-graph weighted by
the number of digit-extensions, and $\E[x^{K}]\sim\beta_{\mathrm{sat}}(x)^n$ by
the same transfer-operator computation; $\beta_{\mathrm{sat}}$ is analytic and
$\beta_{\mathrm{sat}}(1)=\lambda^{\mathrm{sat}}$ is simple, so Hwang's
quasi-power theorem gives the stated Gaussian limit. Primitivity is exactly what
preserves the simplicity of the dominant eigenvalue and hence the analyticity of
$\beta_{\mathrm{sat}}$; if the gcd condition fails (e.g.\ only $c_2\ne0$), the
matrix is imprimitive of period $d=\gcd\{i:c_i>0\}$, the digit string splits into
$d$ interleaved independent strings, and the CLT holds for each block with the
same constants---so a CLT still holds, with the variance constant multiplied by
the appropriate $1/d$ factor coming from the block decomposition.
\end{proof}

\subsection{Relaxation~II (non-existence): a forbidden digit value}
\begin{definition}\label{def:forb}
Fix $1\le r\le c_1-1$ (possible only if $c_1\ge2$). The \emph{$r$-restricted}
rule keeps \eqref{eq:legdef} but additionally forbids any digit equal to $r$.
\end{definition}

\begin{proposition}[Non-existence and density under Relaxation~II]
\label{prop:noexist}
Under the $r$-restricted rule the set of representable integers is exactly the
image of the legal strings avoiding the digit $r$; its transfer matrix is
$T^{(r)}=Z'+N'$ where $Z',N'$ delete all extensions by the digit $r$. Then
$T^{(r)}$ is non-negative with Perron root $\lambda^{(r)}<\lambda$, and the number
of representable integers in $[0,H_{n+1})$ is
$\Theta\big((\lambda^{(r)})^{n}\big)=o(H_{n+1})$. Hence the representable
integers have density $0$, so a positive proportion of integers have no
$r$-restricted decomposition. Conditioned on representability (equivalently, on
the digit string drawn from the $T^{(r)}$-subshift), the number of summands is
asymptotically Gaussian with constants $A^{(r)}=g_r'(1)$,
$C^{(r)}=g_r''(1)+g_r'(1)$, $g_r=\log\beta_r$, $\beta_r(x)$ the Perron root of
$Z'+xN'$, provided $T^{(r)}$ is primitive (which holds whenever the remaining
allowed nonzero-digit positions still have index-gcd $1$).
\end{proposition}

\begin{proof}
Deleting a digit value removes edges from the legality graph, so
$T^{(r)}\le T$ with strict inequality; Perron monotonicity gives
$\lambda^{(r)}<\lambda$ (strict because the deleted edge lies on the recurrent
class through $\mathbf 0$, as $r\le c_1-1$ is a free leading digit). The count of
representable integers equals the number of length-$n$ legal $r$-avoiding
strings, which is $\one^\top (T^{(r)})^n\one_{\mathbf 0}=\Theta((\lambda^{(r)})^n)$
by Perron--Frobenius; dividing by $H_{n+1}=\Theta(\lambda^n)$ gives density
$(\lambda^{(r)}/\lambda)^n\to0$. The conditional CLT is the quasi-power theorem
applied to $Z'+xN'$ exactly as in Proposition~\ref{prop:cltI}; primitivity again
guarantees simplicity of $\beta_r(1)=\lambda^{(r)}$ and the analyticity needed
for the quasi-power method. If primitivity fails the period-$d$ block argument of
Proposition~\ref{prop:cltI} again yields a (blockwise) CLT.
\end{proof}

\begin{remark}[Dichotomy]
Combining the three regimes: relaxing legality \emph{upward} (saturated rule,
Def.~\ref{def:sat}) raises the Perron root and destroys uniqueness but preserves
existence; restricting it \emph{downward} (forbidden digit, Def.~\ref{def:forb})
lowers the Perron root, preserves uniqueness on the representable set, but
destroys existence on a density-$1$ set. The rigid rule is the unique threshold
between these: it is the maximal rule for which every integer is representable
and the minimal rule for which representations are unique. In all three regimes
the number of summands satisfies a Gaussian CLT with mean and variance growing
linearly in $n$, the growth constants being logarithmic derivatives of the Perron
root of the corresponding tilted transfer matrix.
\end{remark}

\section{Part (b): the limiting gap distribution}

\subsection{The stationary (Parry) measure}
By Lemma~\ref{lem:bij}, choosing $N$ uniformly in $[H_n,H_{n+1})$ is the same as
choosing a uniform legal string of length $n$. As $n\to\infty$, the empirical law
of the digit string near a typical position converges to the
\emph{measure of maximal entropy} (Parry measure) of the SFT. We record it in the
``from$\to$to'' convention. Let $A=T^\top$, so $A_{s,s'}$ is the number of legal
one-digit extensions leading from state $s$ to state $s'$; let $r$ be the right
Perron eigenvector of $A$ (equivalently the left Perron eigenvector of $T$),
$Ar=\lambda r$, normalized $r_{\mathbf 0}=1$, and let $\ell$ be the left Perron
eigenvector of $A$ normalized by $\ell^\top r=1$. Then the digit process is the
stationary Markov chain on $\mathcal S$ with transition matrix and stationary
distribution
\begin{equation}\label{eq:parry}
P_{s\to s'}=\frac{1}{\lambda}\,A_{s,s'}\,\frac{r_{s'}}{r_s},\qquad
\pi_s=\ell_s\, r_s\ \ (\textstyle\sum_s\pi_s=1).
\end{equation}
Each row of $P$ sums to $1$: $\sum_{s'}P_{s\to s'}=\frac1{\lambda r_s}\sum_{s'}A_{s,s'}r_{s'}=\frac{(Ar)_s}{\lambda r_s}=1$;
and $\pi P=\pi$ because $\ell^\top A=\lambda\ell^\top$. (That \eqref{eq:parry} is
the limit of the uniform measure on length-$n$ paths is the Shannon--Parry
theorem for irreducible SFTs; primitivity from Lemma~\ref{lem:PF} supplies
irreducibility and aperiodicity.) Writing $A=A_0+A_+$ where $A_0=Z^\top$ records
$0$-digit extensions and $A_+=N^\top$ records nonzero ones, split $P=Q+B$:
\begin{equation}\label{eq:QB}
Q_{s\to s'}=\frac{1}{\lambda}\,(A_0)_{s,s'}\,\frac{r_{s'}}{r_s},\qquad
B_{s\to s'}=\frac{1}{\lambda}\,(A_+)_{s,s'}\,\frac{r_{s'}}{r_s}.
\end{equation}

\subsection{The gap distribution via renewal}
A \emph{summand position} is an index $i$ with $a_i\ge1$; the gap between two
\emph{consecutive} summand positions $i>j$ is $i-j$. (When $a_i\ge2$ the index $i$
hosts several summands; consecutive ones are at gap $0$. The problem's pairing is
of consecutive distinct summand \emph{indices}, the convention under which the
Fibonacci answer $(\varphi-1)^2\varphi^{-(k-2)}$ holds; we therefore count gaps
between distinct nonzero positions, $k\ge1$.) Let $\mathbf 1_B=B\one$ be the
column of nonzero-emission probabilities, and let
\begin{equation}\label{eq:after}
\boldsymbol\beta=\pi B/q,\qquad q=\pi\,\mathbf 1_B,
\end{equation}
so $q$ is the stationary density of nonzero digits and $\boldsymbol\beta$ is the
distribution of the chain's state \emph{immediately after} emitting a nonzero
digit.

\begin{theorem}[Limiting gap distribution]\label{thm:gap}
For every admissible $(c_1,\dots,c_L)$ with $\gcd\{i:c_i>0\}=1$ and every
$k\ge1$,
\begin{equation}\label{eq:Pk}
P(k)=\lim_{n\to\infty}P_n(k)
=\boldsymbol\beta\,Q^{\,k-1}\,\mathbf 1_B ,
\end{equation}
with $q,\boldsymbol\beta,Q,\mathbf 1_B$ from \eqref{eq:QB}--\eqref{eq:after}.
Moreover:
\begin{enumerate}
\item[(i)] $\sum_{k\ge1}P(k)=1$.
\item[(ii)] (Geometric tail / closed form for $k\ge L$.) Since $Z$ has Perron
root $1$ (Lemma~\ref{lem:PF}), the matrix $Q$ has Perron root $1/\lambda$, and
for all $k\ge L$ the all-zero state dominates so that
\begin{equation}\label{eq:tail}
P(k)=C\,\lambda^{-k},\qquad k\ge L,\qquad
C=\lambda^{L}\,\boldsymbol\beta\,Q^{\,L-1}\,\mathbf 1_B ,
\end{equation}
a closed-form rational expression in the entries of $u,v$ and hence in
$(c_1,\dots,c_L)$. For $1\le k<L$ the value \eqref{eq:Pk} is the same finite
spectral expression but is \emph{not} purely geometric; it depends on which
$c_i$ vanish through the entries of $Q,B$.
\item[(iii)] (Fibonacci.) For $L=2$, $c_1=c_2=1$ (so $\lambda=\varphi$),
\eqref{eq:Pk} gives $P(1)=0$ and $P(k)=\varphi^{-k}$ for $k\ge2$; since
$\varphi^{-2}=(\varphi-1)^2$, this is exactly $P(k)=(\varphi-1)^2\varphi^{-(k-2)}$.
\end{enumerate}
\end{theorem}

\begin{proof}
\emph{Formula \eqref{eq:Pk}.} Fix $k\ge1$. In the stationary chain, the event
``a nonzero digit at position $0$, then $0$'s at positions $1,\dots,k-1$, then a
nonzero digit at position $k$'' (positions read in the descending-index
direction of the decomposition) has probability, by the Markov property and
\eqref{eq:QB},
\[
N(k)=\pi\,B\,Q^{\,k-1}\,\mathbf 1_B
=q\,\boldsymbol\beta\,Q^{\,k-1}\,\mathbf 1_B .
\]
Here the first factor $\pi B$ accounts for the nonzero digit at position $0$ and
the landing state; $Q^{\,k-1}$ for the $k-1$ intervening zeros; $\mathbf 1_B$ for
the terminating nonzero digit. By the ergodic theorem for the irreducible
aperiodic chain, along a uniform length-$n$ string the empirical fraction of
nonzero positions whose next nonzero neighbor is at distance exactly $k$ converges
a.s.\ to $N(k)/q$ (the normalization $q$ being the density of nonzero positions,
i.e.\ of ``gap left-endpoints''). Averaging over $N\in[H_n,H_{n+1})$, which is the
uniform string measure, and using bounded convergence, gives
$P_n(k)\to N(k)/q$, which is \eqref{eq:Pk}. The convergence of $P_n$ to the
stationary functional, and the replacement of the finite-$n$ string law by the
Parry measure, are justified because the boundary effects (the top $O(1)$ and
bottom $O(1)$ digits) contribute $O(1/n)$ to the gap counts.

\emph{(i)} Conditioned on position $0$ being nonzero (probability $q$ under
$\pi$), the next nonzero position is at some finite distance a.s.\ (irreducibility
ensures nonzero digits recur), so $\sum_{k\ge1}N(k)=q$, i.e.\
$\sum_{k\ge1}P(k)=\sum_{k\ge1}\boldsymbol\beta Q^{k-1}\mathbf 1_B
=\boldsymbol\beta(I-Q)^{-1}\mathbf 1_B=1$. Here $\sum_{k\ge1}Q^{k-1}=(I-Q)^{-1}$
converges since $Q$ has spectral radius $1/\lambda<1$, and the value $1$ follows
from $(Q+B)\one=\one$ (each row of the stochastic matrix $P=Q+B$ sums to $1$),
which gives $\mathbf 1_B=B\one=(I-Q)\one$, hence
$\boldsymbol\beta(I-Q)^{-1}\mathbf 1_B=\boldsymbol\beta\one=1$ because
$\boldsymbol\beta=\pi B/q$ is a probability vector
($\boldsymbol\beta\one=\pi B\one/q=\pi\mathbf 1_B/q=q/q=1$).

\emph{(ii)} By Lemma~\ref{lem:PF} the only recurrent state of $Z$ (hence of $Q$)
is $\mathbf 0$, reached from any state by at most $L-1$ forced zero-steps. Thus
for $k-1\ge L-1$, i.e.\ $k\ge L$, the vector $\boldsymbol\beta Q^{\,k-1}$ is
proportional to the left Perron vector of $Q$ supported on $\mathbf 0$, and
$Q$ acts on it by the scalar $1/\lambda$ at each further step:
$\boldsymbol\beta Q^{\,k}=\lambda^{-1}\boldsymbol\beta Q^{\,k-1}$ for $k\ge L$.
Hence $P(k+1)=\lambda^{-1}P(k)$ for $k\ge L$, giving \eqref{eq:tail} with the
stated $C$. For $1\le k<L$ the chain has not yet collapsed to $\mathbf 0$ and
\eqref{eq:Pk} retains the transient (non-geometric) contributions of the other
eigenvalues of $Q$ (all $0$, but entering through the nilpotent part), whose
size is governed by which $Z_{s',s}$ (equivalently which $c_i$) vanish.

\emph{(iii)} For $L=2,c_1=c_2=1$: $\mathcal S=\{(0),(1)\}$, and the legal windows
are $(0,0),(0,1),(1,0)$ (the window $(1,1)$ is illegal). Ordering the states
$(0),(1)$ with columns = ``from'' and rows = ``to'',
\[
Z=\begin{pmatrix}1&1\\0&0\end{pmatrix},\quad
N=\begin{pmatrix}0&0\\1&0\end{pmatrix},\quad
T=\begin{pmatrix}1&1\\1&0\end{pmatrix}.
\]
Here $T$ is symmetric, so $A=T^\top=T$ and the right Perron eigenvector of $A$ is
$r=(1,\varphi-1)^\top$ (normalized $r_{(0)}=1$), with $\lambda=\varphi$. The left
Perron eigenvector $\ell\propto(\varphi,1)$ normalized by $\ell^\top r=1$ gives the
stationary marginal $\pi=(\pi_{(0)},\pi_{(1)})=\big(\tfrac{5+\sqrt5}{10},\tfrac{5-\sqrt5}{10}\big)$.
We need only the Parry matrices \eqref{eq:QB} (with $A_0=Z^\top$, $A_+=N^\top$):
\[
Q_{s\to s'}=\tfrac1\varphi (A_0)_{s,s'}\tfrac{r_{s'}}{r_s}
=\begin{pmatrix}\varphi^{-1}&0\\1&0\end{pmatrix},\qquad
B_{s\to s'}=\tfrac1\varphi (A_+)_{s,s'}\tfrac{r_{s'}}{r_s}
=\begin{pmatrix}0&\varphi^{-2}\\0&0\end{pmatrix},
\]
the entry $B_{(0)\to(1)}=\frac1\varphi\cdot\frac{r_{(1)}}{r_{(0)}}=\frac{\varphi-1}{\varphi}=\varphi^{-2}$.
Hence $\mathbf 1_B=B\one=(\varphi^{-2},0)^\top$, and after one nonzero emission the
chain is in state $(1)$, so $\boldsymbol\beta=\pi B/q=(0,1)$. Therefore, using
$P(k)=\boldsymbol\beta Q^{\,k-1}\mathbf 1_B$: for $k=1$,
$\boldsymbol\beta\,\mathbf 1_B=(0,1)\cdot(\varphi^{-2},0)^\top=0$, so $P(1)=0$;
for $k\ge2$, $\boldsymbol\beta Q=(1,0)$ and $(1,0)Q^{\,k-2}=(\varphi^{-(k-2)},0)$,
whence
\[
P(k)=(\varphi^{-(k-2)},0)\cdot(\varphi^{-2},0)^\top=\varphi^{-(k-2)}\varphi^{-2}
=\varphi^{-k}\qquad(k\ge2).
\]
This matches \eqref{eq:tail} with $C=\lambda^{2}P(2)=\varphi^{2}\varphi^{-2}=1$.
Finally $\varphi^{-2}=(\varphi-1)^2$ because $\varphi^{-1}=\varphi-1$, so
$P(k)=(\varphi-1)^2\varphi^{-(k-2)}$ for $k\ge2$, the known Fibonacci gap law.
\end{proof}

\begin{remark}[Numerical confirmation]
Direct enumeration of $[H_n,H_{n+1})$ and the spectral formula \eqref{eq:Pk}
agree, and exhibit the exact geometric tail $P(k)=C\lambda^{-k}$ for $k\ge L$:
e.g.\ $C=1$ for $(1,1)$, $C=2$ for $(2,1)$, $C\approx1.156$ for $(1,1,1)$, with in
all cases ratio $P(k+1)/P(k)=1/\lambda$ for $k\ge L$. For the base-$b$ case
$L=1,c_1=b$ one gets $P(k)=\tfrac{b-1}{b}\,b^{-(k-1)}$ for all $k\ge1$ (here
$\lambda=b$, $C=b-1$), the digit-independence answer.
\end{remark}

\section{Part (c): the limiting longest-gap distribution}

\subsection{Reduction to extreme values of weakly dependent gaps}
Fix $N$ uniform in $[H_n,H_{n+1})$. Its decomposition has $K(N)$ summands, hence
$K(N)-1$ gaps $G_1,\dots,G_{K-1}$ (between consecutive nonzero positions). By
\eqref{eq:ACrigid}, $K(N)=An+O(\sqrt n)$ with high probability, where $A=q>0$;
thus the number of gaps is $m_n:=An(1+o(1))$. The longest gap is
$M_n=\max_{1\le t\le K-1}G_t$. The gaps form a stationary, exponentially
$\phi$-mixing sequence: indeed they are a deterministic functional of the Parry
Markov chain \eqref{eq:parry}, whose transition matrix $P$ has second eigenvalue
modulus $\rho<1$ (Perron simplicity, Lemma~\ref{lem:PF}), so for the chain
$\phi(t)\le K_0\rho^{\,t}$ for constants $K_0,\rho$; the gap between two nonzero
positions $\ge \tau$ apart depends on chain-segments separated by $\ge\tau-L$
steps, giving the same geometric mixing for $\{G_t\}$.

\subsection{Tail asymptotics and the Gumbel limit}
By Theorem~\ref{thm:gap}(ii), the marginal tail is exactly geometric for
$k\ge L$:
\begin{equation}\label{eq:tailprob}
\Pr(G_t\ge k)=\sum_{j\ge k}P(j)=C\,\frac{\lambda^{-k}}{1-\lambda^{-1}}
=\frac{C}{\lambda-1}\,\lambda^{-(k-1)}\qquad(k\ge L).
\end{equation}
Set $\theta:=\frac{C}{\lambda-1}$, so $\Pr(G\ge k)=\theta\,\lambda^{-(k-1)}$ for
$k\ge L$.

\begin{theorem}[Longest gap: doubly-exponential limit]\label{thm:longest}
With $A=q$ the summand density and $\theta=C/(\lambda-1)$ as above, define the
centering
\begin{equation}\label{eq:bn}
b_n=\frac{\log\!\big(\theta A n\big)}{\log\lambda}=\log_\lambda(\theta A n).
\end{equation}
Then for every fixed $y\in\R$,
\begin{equation}\label{eq:gumbel}
\lim_{n\to\infty}\Pr\!\Big(M_n\le b_n+\frac{y}{\log\lambda}\Big)
=\exp\!\big(-\lambda^{-\lfloor y\rfloor-\{b_n\}\,\mathrm{adj}}\big)
\xrightarrow[\text{along subsequences fixing }\{b_n\}]{}
\exp\!\big(-\lambda^{-y}\big),
\end{equation}
i.e.\ $M_n-b_n$ converges (along subsequences where the fractional part of $b_n$
stabilizes, as is unavoidable for integer-valued maxima with geometric tails) to
a Gumbel-type law: for integer thresholds $k=b_n+y$,
$\Pr(M_n\le k)=\exp(-\theta A n\,\lambda^{-k})(1+o(1))$. In particular
$\E[M_n]=\log_\lambda(\theta A n)+\frac{\gamma}{\log\lambda}+o(1)$ up to the
standard $O(1)$ oscillation, where $\gamma$ is the Euler--Mascheroni constant;
thus the longest gap grows like $\log n/\log\lambda$.
\end{theorem}

\begin{proof}
\emph{Poisson approximation (Chen--Stein).} For a threshold $k=k_n\to\infty$, let
$W=\#\{t:G_t\ge k\}=\sum_{t=1}^{m_n}\one\{G_t\ge k\}$. Then $\{M_n<k\}=\{W=0\}$.
By stationarity and \eqref{eq:tailprob},
$\E[W]=m_n\Pr(G\ge k)=An(1+o(1))\cdot\theta\lambda^{-(k-1)}$. Choose $k=b_n+y$
with $b_n$ as in \eqref{eq:bn}, so that $\E[W]\to \theta A\cdot\lambda^{-(y-1)}\,$
$\cdot\,$(absorbing constants) $\to\Lambda(y):=\lambda^{-y+1}\cdot\frac{\theta A}{\theta A}$;
more precisely $\E[W]= \theta A n\,\lambda^{-(k-1)}(1+o(1))=\lambda^{-(k-b_n)+0}(1+o(1))
=\lambda^{-(y-?)}$, and writing $k=b_n+y$ with $b_n=\log_\lambda(\theta A n)$ gives
$\E[W]=\lambda^{\,1-y}(1+o(1))$. The Chen--Stein bound for sums of indicators of a
$\phi$-mixing sequence (Arratia--Goldstein--Gordon) gives
\[
\big|\Pr(W=0)-e^{-\E W}\big|\le b_1+b_2+b_3,
\]
where $b_1=\sum_{t}\sum_{|t'-t|\le d}\Pr(G_t\ge k)\Pr(G_{t'}\ge k)$,
$b_2=\sum_t\sum_{0<|t'-t|\le d}\Pr(G_t\ge k,G_{t'}\ge k)$ measure local
clustering within a window of size $d$, and
$b_3=\sum_t\E\big|\Pr(G_t\ge k\mid \mathcal F_{>t+d})-\Pr(G_t\ge k)\big|$
measures long-range dependence. Take $d=d_n\to\infty$ with $d_n=o(n)$ and
$\rho^{d_n}=o(1/n)$ (possible since $\rho<1$, e.g.\ $d_n=\lceil 2\log n/\log(1/\rho)\rceil$).
Then $b_3\le m_n K_0\rho^{d_n}=o(1)$ by geometric mixing. For $b_1$:
$b_1\le m_n\,d_n\,\Pr(G\ge k)^2=O(n\,d_n\,\lambda^{-2k})=O(d_n\,n^{-1})=o(1)$ since
$\lambda^{-k}=\Theta(1/n)$. For $b_2$ (clustering): two gaps both $\ge k$ within
$d_n$ of each other require two long zero-runs nearby; by the Markov structure
$\Pr(G_t\ge k,G_{t'}\ge k)\le \Pr(G\ge k)\cdot\max_s(Q^{\,k-1})_{s\,\mathbf0}
\le \Pr(G\ge k)\cdot C'\lambda^{-(k-L)}$, so
$b_2\le m_n d_n \Pr(G\ge k)C'\lambda^{-(k-L)}=O(n\,d_n\,\lambda^{-2k+L})=o(1)$ by
the same accounting. Hence $\Pr(M_n<k)=\Pr(W=0)=e^{-\E W}(1+o(1))
=\exp\!\big(-\theta A n\,\lambda^{-(k-1)}\big)(1+o(1))$.

\emph{Conclusion.} Substituting $k=b_n+y$ ($y\in\Z$, since $M_n\in\Z$) gives
$\Pr(M_n\le b_n+y)=\exp\!\big(-\lambda^{1-y}\cdot\lambda^{-\{b_n\}}\big)(1+o(1))$,
the discrete double-exponential (Gumbel) law; the unavoidable dependence on the
fractional part $\{b_n\}$ is the classical lattice oscillation for maxima of
integer geometric variables. The mean expansion
$\E[M_n]=\log_\lambda(\theta A n)+\gamma/\log\lambda+o(1)$ follows by summation
$\E[M_n]=\sum_{k\ge0}\Pr(M_n>k)$ and the Gumbel mean, again up to the $O(1)$
oscillating term. Thus $M_n=\dfrac{\log n}{\log\lambda}+O(1)$, with the fine
fluctuations governed by the Gumbel law of parameter
$\theta=\dfrac{C}{\lambda-1}$, where $C$ is the gap prefactor \eqref{eq:tail}.
\end{proof}

\begin{remark}
All constants in Part~(c) are explicit functions of $(c_1,\dots,c_L)$: $\lambda$
is the Perron root of \eqref{eq:charpoly}; $A=q=\pi\mathbf1_B=g'(1)$ is the
summand density from \eqref{eq:ACrigid}; and $C$ (hence
$\theta=C/(\lambda-1)$) is the gap prefactor \eqref{eq:tail}, a rational function
of the Perron eigenvector entries $u,v$ of the transfer matrix $T$. For the
Fibonacci case $\lambda=\varphi$, $C=1$, $A=q=1-\varphi^{-2}=2-\varphi$,
$\theta=1/(\varphi-1)=\varphi$, recovering the known
$M_n\sim\log_\varphi n$ Gumbel behavior.
\end{remark}

\end{document}
